% !TeX root = ../../arithmetic-and-computation.tex
\section{How to use this supplement}
This supplement contains the detailed packet arithmetic, historical
exclusion records, auxiliary lattice arguments and execution provenance
used by Papers I and II and the supporting Hodge note. It is organized by
mathematical role. The historical partial computations are explicitly
separated from the final complete packet theorem in Paper I.

Notation in packet sections is $M=\mathbb Q(\zeta_{42})$, $L=M^+$,
$t=\zeta_{42}+\zeta_{42}^{-1}$ and
$\pi=t^4+t^3-4t^2-2t+5$ of norm $43$. In quotient-family material,
$\pi$ instead denotes $1-z^{14}$, the generator over $3$ specified in
Paper II; these symbols must not be identified. Embeddings are ordered
$1,5,11,13,17,19$; $K=\mathbb Q(\sqrt{-3})$.

The roles of evidence are distinct: a finite-domain proof establishes
exhaustiveness; exact arithmetic or outward enclosures justify decisions;
source/runtime and retained-output records describe the execution to which
the claim is bound. A hash establishes file identity, not mathematical
correctness or absence of a historical transient replacement.

\subsection{Result-to-evidence guide}
\begin{description}[leftmargin=1em]
\item[Packet domain.] Paper I, Theorem~\externalref{packet}{thm:fullpacketarithmetic},
uses the order and module calculations below and the full-packet arithmetic
sources in Section~\ref{app:computer}.
\item[Curve intersection.] Paper I, Theorem~\externalref{packet}{thm:fullpacketgeometry},
uses its two-region lemma, positive vertex realizations and both complete
negative populations. The final receipt is
\path{experiments/continuous_v83/whole_packet_assessment/final-adoption.json}.
The historical execution premise in Paper I remains required.
\item[Principal geometry.] The common-lattice, stabilizer and map proofs
are in Paper I. Their named exact regressions and the supplementary
individual-lattice arguments below do not replace those proofs.
\item[Finite quotients and descent.] Paper II names the adopted exact
finite tables within its proofs. Subgroup, underlying, polarized, orbit
and field counts are different assertions, with different proof inputs.
\item[Jacobian exclusion.] Paper II, Theorem~\externalref{quotients}{pp:thm:jacobiancomplete},
retains its model, theta-enclosure and execution contract. Counts cannot be
pooled across the four models. The $156$ smaller-order cases remain open.
\item[Hodge tensors.] The Hodge note gives the algebraic constructions and
integral relation arguments; source identity alone does not prove either.
\end{description}
This is a navigation guide, not a newly completed transitive evidence
archive. The repository remains private. Historical source paths are
relative to its root; a release must provide the complete dependencies
and accurate access instructions before claiming reproducibility.
